\documentclass[11pt,a4paper]{article}
\usepackage[T1]{fontenc}
\usepackage{lmodern,amsmath,amssymb,amsthm,mathtools,booktabs,microtype}
\usepackage[margin=26mm]{geometry}
\usepackage[colorlinks=true,linkcolor=blue,urlcolor=blue]{hyperref}
\setlength{\emergencystretch}{2em}
\newtheorem{theorem}{Theorem}[section]
\newtheorem{lemma}[theorem]{Lemma}
\newtheorem{corollary}[theorem]{Corollary}
\theoremstyle{remark}\newtheorem{remark}[theorem]{Remark}
\newcommand{\Q}{\mathbb Q}\newcommand{\Z}{\mathbb Z}
\title{Local obstructions to shared elliptic reference tests}
\author{}\date{16 November 2026}
\begin{document}\maketitle
\begin{abstract}
We give exact local obstructions to two detector factors using the
shared inverse-reference function of a three-coordinate elliptic chart.
In the squared-point factor, two negative-valuation summands cancel
while the factor equals one and its own ratio comparison is exact.
In the odd-cube factor, Hensel's lemma produces whole-factor zeros
and nonzero values of arbitrarily large valuation, even with that
factor's own comparison. The examples use actual elliptic points and
actual odd contact weights at a good split prime above the prescribed
cutoff. They disprove the corresponding factor-level noncancellation
assertions. They do not impose the global cyclic-index, conjugate-factor,
height, or full pole constraints of a defining-polynomial source.
\end{abstract}

\section{Elliptic and contact data}
Put
\[
 K=\Q(\rho),\quad\rho^2=2,\quad D=11-6\rho,
 \quad E:y^2=x^3-Dx,
\]
\[
 P=(a,b)=\left(\frac{57}{8}-\frac72\rho,
                     -\frac{165}{8}+\frac{417}{32}\rho\right),
 \quad c=b/a,\quad\delta=D/a,\quad T_0=(0,0).
\]
The complete reference chart~\cite{Chart} has coordinates
$S,T_H,v$ on $E\setminus\{O,T_0,P,P+T_0\}$, where
\[
 S=\frac{y(H)}{x(H)},\qquad T_H=\frac{y(H-P)}{x(H-P)},
 \qquad v=\frac{y(H)+b}{x(H)-a}+c.
\]
The last expression is interpreted by its regular tangent extension
at $H=-P$. The chart equations imply that
\[
                 \ell=v^{-1}=-\frac{v+T_H-S}{\delta}
\]
is linear in its coordinates.

Fix the rational prime and residue data
\begin{equation}\label{eq:prime}
 p=76537,\qquad \rho\equiv26935\ \hbox{or}\ 49602\pmod p,
 \qquad i\equiv8042\pmod p.
\end{equation}
The integer $p$ is prime, $p>76176$, and the displayed residues satisfy
$\rho^2=2$ and $i^2=-1$ modulo $p$. Each lifts uniquely to a root
with that residue in $\Q_p$. Both resulting embeddings of $K$
have $D,a,b,\delta$ units, and $E$ has good reduction.
Write $\operatorname{ord}$ for the normalized $p$-adic valuation.

Use the actual contact weight
\begin{equation}\label{eq:contact}
 F(t)=\prod_{\theta\in\{-5,-2,-1,0,2,3,4,7\}}(t-\theta),
 \qquad w(t)=\frac{F(t)}{(t-1)^6}.
\end{equation}
It satisfies $w(0)=0$ and $w'(0)=1680$, a $p$-adic unit.
Consequently every sufficiently small nonzero $w\in\Q_p$ is
$w(t)$ for a unique $t\in p\Z_p$, with
\begin{equation}\label{eq:contactorder}
                          \operatorname{ord}(t)=\operatorname{ord}(w).
\end{equation}
This follows either from the local inverse theorem or directly by
Hensel lifting $F(t)-w(t-1)^6$ from $t=0$.

\section{Exact local points and ratios}
For a numerator point use
\[
 f=1/x(U),\qquad h=-y(U)/x(U)^2,
 \qquad h^2=f-Df^3,
\]
and put $z_U=f/h=-x(U)/y(U)$ when $U\ne O$ is formal.

\begin{lemma}\label{lem:local}
For every nonzero $z\in p\Z_p$ there is a unique formal point
with parameter $z$, whose reciprocal coordinates satisfy
\[
 f+D z^2f^2=z^2,\qquad h=f/z,\qquad f/z^2\in1+p\Z_p.
\]
In particular $f$ is a square in $\Q_p$ and
$\operatorname{ord}(f)=2\operatorname{ord}(z)$.

Take a formal point $Q$ of parameter $z_Q$ and put $H=T_0+Q$.
Then $H$ lies in the full reference-chart domain, and
\begin{equation}\label{eq:localref}
 S=1/z_Q,\qquad
 v=\frac{\delta(h_Q+c f_Q)}{1+\delta f_Q},\qquad
 \operatorname{ord}(\ell)=-\operatorname{ord}(z_Q).
\end{equation}
For any nonidentity formal numerator $U$, the two ratio rows
\[
 hL+fS=0,\qquad(1-Df^2)L+hS=0
\]
hold with the exact value
\begin{equation}\label{eq:rawratio}
                              L=-z_U/z_Q.
\end{equation}
\end{lemma}

\begin{proof}
The equation $f+D z^2f^2-z^2=0$ has derivative one modulo $p$ at
$f=0$, so Hensel's lemma gives its unique small root. It is
$z^2$ times a unit congruent to one modulo $p$. Substitution of
$h=f/z$ gives the reciprocal curve equation. Since $p$ is odd,
that unit has a square root, proving the square assertion.

Translation by $T_0$ is
\[
 (x,y)\longmapsto(-D/x,Dy/x^2).
\]
Thus $x(H)=-Df_Q$, $y(H)=-Dh_Q$, and $S=h_Q/f_Q=1/z_Q$.
The formula for $v$ follows by substitution in its definition.
Its leading term is $\delta z_Q$ times a unit, proving the order.
The point $H$ reduces to $T_0$ but is not equal to it; the other
excluded reference points have distinct reductions at the chosen
prime. Finally the first ratio row gives~\eqref{eq:rawratio},
and the curve equation gives the second.
\end{proof}

\section{Cancellation in the squared-point factor}
Consider
\begin{equation}\label{eq:squared}
           \mathcal T=1+q^2C^4+q\ell^2+q^3f^2,
           \qquad q=w^{-1},\qquad C=L-\beta.
\end{equation}

\begin{theorem}\label{thm:squared}
Let $\beta\in\Z_p^\times$ and let $k$ be any positive odd integer.
There are local points $H,U$ satisfying all reference and ratio
conditions of Lemma~\ref{lem:local}, together with an actual
contact parameter $t$, for which
\[
 \operatorname{ord}(t)=3k,\qquad L=\beta,\qquad
 \operatorname{ord}(\ell)=-k,\qquad\operatorname{ord}(f)=2k,
\]
but
\[
 \operatorname{ord}(q\ell^2)=\operatorname{ord}(q^3f^2)=-5k,
 \qquad q\ell^2+q^3f^2=0,\qquad\mathcal T=1.
\]
\end{theorem}

\begin{proof}
Choose $z_Q$ of order $k$ and $z_U=-\beta z_Q$.
Lemma~\ref{lem:local} gives $L=\beta$ and hence $C=0$.
Set
\[
                              q=i\ell/f.
\]
Then
\[
                q\ell^2+q^3f^2
                 =\frac{i\ell^3}{f}(1+i^2)=0.
\]
The orders are as stated, and $w=q^{-1}$ has positive odd order
$3k$. Equations~\eqref{eq:contact}--\eqref{eq:contactorder}
realize it by the required contact parameter.
\end{proof}

The factor has valuation zero although its shared-reference term has
large negative valuation. Thus its valuation cannot be used to bound
the valuations of its individual summands. The comparison here is
already exact; its failure is not caused by choosing a different
comparison expression in another factor.

\section{Whole-factor obstruction for the odd-cube test}
For $s\in\{1,3\}$ consider
\begin{equation}\label{eq:cube}
      \mathcal A_s=w^s+C^2+w^{-1}\ell^2+w^{-3}f^3,
      \qquad C=L-\beta.
\end{equation}
The identity $f^3=(f-h^2)/D$ allows the point term to be used in
a quadratic clearing with the usual point helpers. The following
obstruction concerns the factor itself, independently of that encoding.

\begin{theorem}\label{thm:cube}
Let $\beta\in\Z_p^\times$. Choose positive integers $k,m$ of
opposite parity. There are local reference and numerator points as
in Lemma~\ref{lem:local}, with
\[
 \operatorname{ord}(\ell)=-k,\quad\operatorname{ord}(f)=2m,
 \quad\operatorname{ord}(L)=m-k,
\]
such that, for each $s=1$ or $3$, an actual contact weight $w$ of
positive odd order $k+3m$ satisfies $\mathcal A_s=0$.
With the points and target fixed, there are also actual contact
weights of that same order for which $\mathcal A_s\ne0$ has
arbitrarily large positive valuation.
\end{theorem}

\begin{proof}
Choose formal parameters of orders $k,m$ and use
Lemma~\ref{lem:local}. Put $c_0=\operatorname{ord}(C)$; then
\[
                         c_0\ge\min(m-k,0).
\]
Choose a square root of $f$ and set
\[
                   w_0=i f^{3/2}/\ell,
                   \qquad\nu=\operatorname{ord}(w_0)=k+3m.
\]
Multiplication of~\eqref{eq:cube} by $w^3$ gives
\[
 P_s(w)=w^{s+3}+C^2w^3+\ell^2w^2+f^3.
\]
The substitution $w=w_0Y$ yields the exact identity
\begin{equation}\label{eq:hensel}
 \frac{P_s(w_0Y)}{f^3}
   =1-Y^2+\epsilon_3Y^3+\epsilon_{s+3}Y^{s+3},
\end{equation}
where
\[
 \epsilon_3=\frac{C^2w_0^3}{f^3},\qquad
 \epsilon_{s+3}=\frac{w_0^{s+3}}{f^3}.
\]
Their orders satisfy
\[
 \begin{split}
 \operatorname{ord}(\epsilon_3)&=3k+3m+2c_0>0,\\
 \operatorname{ord}(\epsilon_4)&=4k+6m>0,\qquad
 \operatorname{ord}(\epsilon_6)=6k+12m>0.
 \end{split}
\]
For the first strict inequality, the lower bound is $k+5m$ if
$m<k$, and $3k+3m$ if $m\ge k$.
The reduction of~\eqref{eq:hensel} is $1-Y^2$.
Its roots $1,-1$ are simple, so each lifts to an exact root.
The resulting $w$ has order $\nu$, which is odd by the parity
assumption, and is realized by~\eqref{eq:contactorder}.

For the last assertion, let $Y_*$ be one lifted root and replace
it by $Y_*+p^N$ with $N$ sufficiently large. The polynomial in
\eqref{eq:hensel} has unit derivative at $Y_*$, so its value has
order exactly $N$. Consequently
\[
                 \operatorname{ord}(\mathcal A_s)
                      =N+6m-3\nu=N-3k-3m.
\]
This is finite and arbitrarily large. The weight still has order
$\nu$, and every such weight is an actual local contact weight.
\end{proof}

For example, $k=1,m=2,\beta=4$ gives a contact of order seven.
Its raw ratio has order one while $C$ is a unit, so the whole
factor can have arbitrarily large valuation despite a comparison
which fails modulo $p$. The conclusion persists for both displayed
baseline powers and after excluding zero values of the factor.

\section{Verification and precise scope}
The standalone checker in
\nolinkurl{negative_results/computation/shared_reference_tests_v43}
proves primality by trial division, checks both roots of $2$ and the
root of $-1$, verifies the elliptic constants, and proves the displayed
symbolic cancellation and normalization identities. It independently
constructs modular formal points in both embeddings. For the
squared-point factor it verifies the cleared cancellation modulo
$p^{55}$ and the contact equation modulo $p^{50}$. For each odd-cube
baseline it verifies the cleared factor modulo $p^{50}$ and the
contact equation modulo $p^{43}$. The exact existence statements
use the simple-root Hensel arguments above, rather than an inference
from finite precision alone.

These results concern one local factor. They do not assert that the
chosen local reference point belongs to the required doubled global
subgroup, or that the local points come from one set of global cyclic
indices. Both conjugate factors, the second detector with its prescribed
argument shift, the trace-height relation, and the full pole witnesses
are not imposed. A zero at one local embedding is not presented as a
nonzero global field element with zero image. Thus the theorems rule
out the specified local noncancellation arguments; they do not prove
unsoundness of a complete global source, obstruct all shared-reference
strategies, or change a defining-polynomial variable bound.

\begin{thebibliography}{9}
\bibitem{Chart}\emph{A three-coordinate quadratic chart for an elliptic
point and its translate}, project manuscript, 15 November 2026,
\nolinkurl{papers/elliptic_charts/three_coordinate_reference_20261115_v40.tex}.
\end{thebibliography}
\end{document}
